\documentclass{article}
\usepackage{float}
% Language setting
\usepackage[german]{babel}
\usepackage[table]{xcolor}
% Set page size and margins
% Replace `letterpaper' with `a4paper' for UK/EU standard size
\usepackage[letterpaper,top=2cm,bottom=2cm,left=3cm,right=3cm,marginparwidth=1.75cm]{geometry}
\usepackage[dvipsnames]{xcolor}
% Useful packages
\usepackage{amsmath}
\usepackage{graphicx}
\usepackage[colorlinks=true, allcolors=blue]{hyperref}
\usepackage{subcaption}
\usepackage{graphicx}
\title{Fragestellungen von Kolloquium Wendisch I.} 
\author{Lil'Jana}
\usepackage{enumitem,amssymb}
\newlist{todolist}{itemize}{2}
\setlist[todolist]{label=$\square$}
\usepackage{pifont}
\newcommand{\cmark}{\ding{51}}%
\newcommand{\xmark}{\ding{55}}%
\newcommand{\done}{\rlap{$\square$}{\raisebox{2pt}{\large\hspace{1pt}\cmark}}%
\hspace{-2.5pt}}
\newcommand{\wontfix}{\rlap{$\square$}{\large\hspace{1pt}\xmark}}


\begin{document}
\maketitle
\subsection*{Wie schreibt man DNA-Sequenzen?}
Man schriebt den Strang bzw. den Hauptstrang immer vom 5' zum 3'-Ende. \newline
Bspw. 
\begin{center}
    5'-AGCTGCA-3' 
\end{center}
\begin{center}
    \textit{(ss-DNA)}
\end{center}
  bzw. \newline
\centerline{5'-AGCTGCA-3'}
\centerline{3'-TCGACGT-5'}
 \begin{center}
    \textit{(ds-DNA)}
\end{center}


\subsection*{Wohin wandert DNA im elektrischen Feld?}
Zur positiven Elektrode (Anode), weil DNA negativ geladen ist.

\subsection*{Welche strukturellen Eigenschaften haben DNA und Polyphosphat gemeinsam? }
\begin{itemize}
    \item Phosphat-Rückgrat
    \item Negative Ladung
    \item Kettenartige Struktur
    \item Konformationen und Flexibilität \textit{(Können in versch. Konformationen vorliegen, weil sie flexibel sind)}
\end{itemize}

\subsection*{Wie färbt man DNA?}
Das ist durch diverse Farbstoffe möglich. Bspw.: 
\begin{itemize}
    \item Ethidiumbromid
    \item SYBR-Green/SYBR-Gold
\end{itemize}

\subsection*{Warum ist G-C stabiler als A-T?}
Wegen der Anzahl der Wasserstoffbrückenbindungen. 
\begin{itemize}
    \item Guanin (G) und Cytosin (C) bilden drei Wasserstoffbrückenbindungen
    \item Adenin (A) und Thymin (T) bilden nur zwei Wasserstoffbrückenbindungen.
\end{itemize}

\subsection*{Wie unterscheidet sich DNA von RNA?}
\begin{itemize}
    \item Zuckerstruktur
    \begin{itemize}
        \item DNA hat Desoxyribose
        \item RNA hat Ribose
    \end{itemize}
    \item Basenzusammensetzung
    \begin{itemize}
        \item DNA hat A, G, C, T
        \item RNA hat A, G, C, U
    \end{itemize}
    \item Struktur
    \begin{itemize}
        \item DNA ist doppelsträngiger Helix
        \item RNA ist einzelsträngig
    \end{itemize}
    \item Funktion
    \begin{itemize}
        \item Zweck der DNA ist langfristige Speicherung und Weitergabe 
        \item Zwecke der RNA sind vielfältiger - u.a. Botenmolekül (mRNA) und Strukturelement und Enzym (rRNA, tRNA)
    \end{itemize}
    \item Stabilität
    \begin{itemize}
        \item DNA ist stabiler
        \item RNA ist instabiler
    \end{itemize}
    \item Vorkommen und Lebensdauer
    \begin{itemize}
        \item DNA ist (bei Eukaryoten) hauptsächlich im Zellkern zu finden und in Mitochondrien und Chloroplasten
        \item RNA ist hauptsächlich im Cytoplasma zu finden und hat eine kürzere Lebensdauer
    \end{itemize}
\end{itemize}

\subsection*{Warum ist DNA stabiler als RNA?}
Aufgrund mehrerer struktureller Unterschiede: 
\begin{itemize}
    \item Desoxyribose vs. Ribose
    \item Doppelsträngigkeit der DNA
    \item Abwesenheit von Uracil in DNA
    \item Abwesenheit von 2'-Hydroxylgruppe in DNA
\end{itemize}

\subsection*{Aufgabe: Die kodierende Region eines Gens weist bezüglich der Basenzusammensetzung 5\% Adenin und 45\% Guanin auf}
Erinnerung: \newline
\centerline{$A+T+G+C=100\%$}
\centerline{$A = T$}
\centerline{$C = G$}
\subsubsection*{1. Was ist der prozentuale Anteil an Cytosin?}
\centerline{$C = G = 45\%$}
\subsubsection*{2. Was ist in der Entsprechenden RNA der prozentuale Anteil an Uracil?} 
\centerline{$A = T\ bzw.\ U = 5\%$}
\subsubsection*{3. Was ist der prozentuale Anteil an Purinen? } 
Purine sind Guanin und Adenin. \newline
\centerline{$A = 5\%$} \newline
\centerline{$A = 45\%$}
\centerline{$45\% + 5\% = 50\%$}

\subsection*{Welcher Phosphatrest wird in die DNA eingebaut?}
Der 5'-Phosphatrest. 

\subsection*{Welcher Phosphatrest wird bei Kinase-Reaktionen transferiert?}
 ATP (Adenosintriphosphat)
 
\subsection*{Was kann man über die Kinetik der DNA-Synthese sagen?}
\begin{itemize}
    \item Semikonservative Replikation
    \item Replikationsgabel
    \item DNA-Polymerasen
    \item 5'-zu-3'-Richtung
    \item Führender vs. verzögerter Strang
    \item bis zu 1.000 Basenpaare pro Sekunde \textit{(in Bakterien)}
\end{itemize} 

\subsection*{DNA-Replikation}
\subsubsection*{1. Ist wie gerichtet?} 
\textbf{5'-zu-3'-Richtung }\newline
Der führende Strang wird kontinuierlich in Richtung der Replikationsgabel synthetisiert. Der verzögerte Strang wird diskontinuierlich synthetisiert. 
\subsubsection*{2. Verhält sich wie bezüglich der Ausgangs-DNA und Tochter-DNA?} 
\textbf{Semikonservativ} \newline
Nach der Replikation besteht jede der beiden neuen DNA-Doppelhelices aus einem alten (mütterlichen) Strang und einem neuen, neu synthetisierten (tochter) Strang.
\subsubsection*{3. Verhält sich wie bezüglich des Ablaufs der Replikation an den beiden DNA-Strängen?} 
Es gibt einen führenden und einen verzögerten Strang.
\begin{table}[h]
    \centering
    \begin{tabular}{|p{7cm}|p{7cm}|} \hline
     führend & verzögert \\ \hline \hline
    Kontinuierliche Synthese in Richtung der Replikationsgabel &  Diskontinuierliche Synthese in Form von Okazaki-Fragmenten, die durch die Replikationsgabel entfernt sind.  \\\hline
\end{tabular}
\end{table} 

\subsection*{Geben Sie die Werte der Parameter der DNA-
Superhelikalitat an? Wie erfolgen die Umwandlungen?}
\centerline{$LK=T+W$}
\begin{figure}[H]
    \centering
    \includegraphics[width=13cm]{bilder/Circular_DNA_Supercoiling.png}
    \label{fig:my_label}
    \caption{https://upload.wikimedia.org/wikipedia/commons/1/1e/Circular_DNA_Supercoiling.png}
\end{figure}
\subsection*{Aufgabe: In einer Probe seien DNA und RNA enthalten. Sie sollen nur die RNA zerstören. Wie gehen Sie vor?}
Man verwendet RNAse. 

\subsection*{Welche(n) Parameter der Superhelikalität der DNA können nur Topoisomerasen verändern?}
\begin{todolist}
\item [\done] Linking Number (LK)
\item Verdrilling Twist (T) 
\item Anzahl der Überkreuzungen Writhe (W)
\end{todolist}
\subsection*{Wie unterscheidet man superhelikale von 'entspannter' DNA?}
\begin{table}[H]
    \centering
    \begin{tabular}{|c||p{5cm}|p{5cm}|} \hline 
         & Superhelikal & Entspannt  \\ \hline\hline
        Verdrillung & Zusätzliche Windungen oder Überdrehungen &  \textbf{keine} Zusätzliche Windungen oder Überdrehungen\\ \hline
        Topologische Eigenschaften & Veränderter Linking Number (LK) & Natürliche Linking Number (LK)\\ \hline
        Funktionelle Bedeutung &  Erleichtert oder behindert je nach Zustand den Zugang zur DNA für Prozesse wie Replikation und Transkription & Repräsentiert den Grundzustand der DNA ohne zusätzliche Spannung \\\hline
    \end{tabular}
    \label{tab:my_label}
\end{table} 
\subsection*{Lückentext}
Bei der Replikation der DNA von E. coli arbeitet die DNA-Polymerase III. sehr genau. Sie synthetisiert nur jedes \textbf{10 hoch 8}te Basenpaar falsch, weil sie selbst über eine Reparaturfunktion verfügt, die \textbf{proofreading} (Korrekturlesefunktion) ist. \newline
Replizierte DNA weist jedoch noch weniger Fehler auf: nur jedes \textbf{10 hoch 9}te Basenpaar falsch, weil es zusätzliche Reparaturmechanismen gibt, z.B. die \textbf{Mismatch-Reparatur}.

\subsection*{Wie können Mutationen induziert werden?}
\subsubsection*{1. Physikalisch, z.B.} 
    Strahlung
    \begin{itemize}
    \item ionisierende
    \item UV
    \item Partikel-
\end{itemize}
\subsubsection*{2. Chemisch, z.B.} 
Verbindungen, die die DNA schädigen oder verändern können, indem sie mit den Basen der DNA interagieren oder die DNA-Struktur direkt beeinflussen. 
\begin{itemize}
    \item Alkylierungsmittel
    \item Basenanaloga
    \item Oxidative Mittel
    \item ...
\end{itemize}
\subsection*{Aufgabe: Ein Protein hätte drei Aminosäuren (1.Methionin, 2. Tryptophan, 3. Tryptophan). Schreiben Sie die Protein- und die entsprechende mRNA-Sequenz auf! Benennen Sie
die vorderen und hinteren Enden! Benennen Sie den Vorgang}
Was ist hier gefordert? 
\begin{enumerate}
    \item Bestimmung der Proteinsequenz 
    \item Bestimmung der Codons durch Ablesen von der Sonne 
    \item Bestimmung der RNA-Sequenz aus den Codons
    \item Benennen wo Ende und Anfang sind
    \item Benennung des Vorgangs
\end{enumerate}
\subsubsection*{1. Bestimmung der Proteinsequenz} 
\centerline{1. Methionin, 2. Tryptophan, 3. Tryptophan }
\subsubsection*{2. Bestimmung der Codons durch Ablesen von der Sonne }  
\centerline{Methionin: AUG} 
\centerline{Tryptophan: UGG} 
\subsubsection*{3. Bestimmung der RNA-Sequenz aus den Codons} 
\centerline{Methionin - Tryptophan - Tryptophan  = AUG - UGG - UGG} 
\subsubsection*{4. Bennenen wo Ende und Anfang sind} 
AUG ist ein Start-Codon, daher ist bei AUG der Anfang und bei dem zweiten UGG das Ende.  
\subsubsection*{5. Benennung des Vorgangs} 
\centerline{Transkription} 

\subsection*{Aufgabe: Bei der Entschlüsselung des genetischen Codes wurden RNA-Moleküle  die entweder nur aus U oder nur aus CU zusammengesetzt waren, translatiert und die entstehenden Proteine charakterisiert. Welche Aminosäure(n) werden durch eine RNA kodiert, die sich durch aufeinander folgende AG's zusammensetzt?}
Was ist hier gefordert?
\begin{enumerate}
    \item Codons bilden 
    \item Codons zu Aminosäuren zuordnen
\end{enumerate} 
\subsubsection*{1. Codons bilden} 
AGAGAGAG...AGAG beinhaltet das Codon AGA. \newline
\subsubsection*{2. Codons zu Aminosäuren zuordnen} 
AGA codiert für Arginin.  

\subsection*{Genetischen Codes: statt einem Triplett codiert nun ein Quadruplett eine Aminosäure. Berechnen Sie wie viele
Quadruplett-Codons möglich waren!}
\centerline{$Anzahl\ der\ Codons = 4^4 = 256 $}

\subsection*{Für ein E. coli-Gen sei mit einer molekular biologischen Technik der Transkriptionsstartpunkt (das erste Nukleotid der mRNA) bestimmt worden und in der angegebenen Sequenz fett markiert. Identifizieren Sie die Promotor-Elemente!}
\centerline{AAAAAATC\underline{TTGACA}CCCAGCATCAGTAGACATAA\underline{TATGAC}CACA\textbf{T}GGACCAAA} 
\noindent Was wird hier gefordert? 
\subsubsection*{Konsensussequenzen der E. coli Promotoren finden}
-10 Region (Pribnow-Box) und -35 Region. \newline
Hier ist die -10 Region: TATGAC \newline
Und die -35 Region: TGACAC

\subsection*{Die RNA-Polymerase ist wesentlich für die Transkription in E. coli. Welchem Enzym der DNA-Replikation (Primase oder DNA-Polymerase III) ähnelt sie mehr?}
Primase und RNA-Polymerase. 
\begin{table}[h] 
    \centering
    \begin{tabular}{|p{3cm}||p{4cm}|p{4cm}|p{4cm}|}\hline
         & Primase &  DNA-Polymerase III & RNA-Polymerase\\ \hline \hline
         Aktivität /was wird synthetisiert & \cellcolor{SpringGreen} RNA-Primer  & neuer DNA-Strang & \cellcolor{SpringGreen} RNA \\\hline
         Vorläufer/Substrate & \cellcolor{SpringGreen} NTPs \textit{(Ribonukleosidtriphosphate)} & dNTPs \textit{(Desoxyribonukleosidtriphosphate)}& \cellcolor{SpringGreen} NTPs \textit{(Ribonukleosidtriphosphate)} \\\hline
         Matrize &\cellcolor{SpringGreen} ssDNA & \cellcolor{SpringGreen} ssDNA & \cellcolor{SpringGreen} ssDNA \\\hline
         Starter oder Primer &\cellcolor{SpringGreen} Kein Primer benötigt & RNA-Primer  & \cellcolor{SpringGreen} Kein Primer benötigt \\\hline
    \end{tabular}
    \label{tab:my_label}
\end{table}

\subsection*{Wie viele verschiedene Aminosäuren gibt es in
Proteinen?}
\textbf{20} verschiedene

\subsection*{Anzahl der während der Proteinbiosynthese (ko-translational) eingebauten Aminosäuren}
\textbf{20 bis 22} verschiedene 

\subsection*{Woraus setzt sich ein Ribosom in Bakterien zusammen?}
\begin{todolist}
\item nur aus RNA 
\item nur aus Protein
\item [\done] aus RNA und Protein
\item [\done] aus 3 RNA Molekülen und aus 1 Protein
\item aus 3 RNA Molekülen und aus vielen Proteinen 
\item [\done] aus 2 Untereinheiten
\end{todolist} 

\subsection*{Ordnen Sie die
Translationsfaktoren EF-G,EF-TU, IF1, IF2, IF3, RF und RRF zeitlich in den Prozess der Translation ein und weisen Ihnen eine Funktion dabei zu. 
}
\begin{enumerate}
    \item Initiationsphase 
    \begin{enumerate}
        \item IF1 (Initiationsfaktor 1)
        \item IF2 (Initiationsfaktor 2)
        \item IF3 (Initiationsfaktor 3)
    \end{enumerate}
    \item Elongationsphase
    \begin{enumerate}
        \item EF-Tu (Elongationsfaktor Thermounstable):
        \item EF-G (Elongationsfaktor G)
    \end{enumerate}
    \item Terminationsphase
    \begin{enumerate}
        \item RF (Release Factor)
        \item RRF (Ribosome Recycling Factor):
    \end{enumerate}
\end{enumerate}
\subsection*{Was ist  die Shine-Dalgarno-Sequenz?
Wo liegt sie? Was ist  ihre Funktion?}
Die Shine-Dalgarno-Sequenz ist eine wichtige purinreiche Signalsequenz, die in der mRNA von Prokaryoten vorkommt. \newline
Die Shine-Dalgarno-Sequenz liegt upstream (also vor) dem Startcodon (AUG) in der mRNA, etwa 4-9 Nukleotide stromaufwärts. Typischerweise befindet sie sich im 5'-untranslatierten Bereich (5'-UTR) der mRNA. \newline
Die Hauptfunktion der Shine-Dalgarno-Sequenz besteht darin, die Bindung der mRNA an das Ribosom zu erleichtern.

\subsection*{Wie sind Ribosomen aus Archaea und Eukaryoten zusammengesetzt?} 
\begin{table}[H]
    \centering
    \begin{tabular}{ |p{2cm}|| p{4cm}| p{4cm}| p{4cm}|}
    \hline
     & Archaea & Bacteria & Eukaryoten \\ \hline \hline
    Typ &  70S & 70S & 80S\\ \hline
    zusammen gesetzt aus.. & 50S + 30S  & 50S + 30S & 60S + 40S  \\\hline
    rRNA & 23S, 5S, 16S & 23S, 5S, 16S & 28S, 5S, 18S  \\\hline
    ribosomale Proteine & 50-60 & 54-58 & 80-100 \\\hline
    Hemmstoffe & Resistenter gegen viele bakterielle Hemmstoffe & Empfindlich gegenüber spezifischen Hemmstoffen wie Chloramphenicol, Streptomycin & Resistenter gegen bakterielle Hemmstoffe, aber empfindlich gegenüber Cycloheximid, Puromycin \\\hline
\end{tabular}
\end{table}

\subsection*{Was ist das katalytisch aktive Prinzip des Ribosoms?}
Das katalytisch aktive Prinzip des Ribosoms ist die ribosomale RNA (rRNA), insbesondere im Peptidyltransferase-Zentrum. Diese rRNA, die keine direkte Proteinbeteiligung zur Katalyse benötigt, sorgt für die Bildung der Peptidbindungen, die die Aminosäuren zu einer wachsenden Polypeptidkette verbinden. Damit fungiert das Ribosom als ein Riboenzym.

\subsection*{Aufgabe: Ordnen Sie die Vorgänge bei der homologen Rekombination. (nicht alle existieren)}
\begin{enumerate}
    \item Doppelstrang-Bruch
    \item Ausbildung von 3'-Überhängen
    \item Eindringen eines Einzelstrangs
    \item Branch Migration
    \item Auflösen der Holliday-Struktur
\end{enumerate}
Es existieren nicht: 
\begin{itemize}
    \item  Eindringen eines Doppelstrangs
    \item Ausbildung von 5'-Überhängen
    \item Wandern der Holliday-Struktur \textit{= (Branch Migration)}
\end{itemize}

\subsection*{Bei der Rekombination zweier homologer DNA-Regionen mit den nach einander liegenden Genen 1, 2, 3 und 4 bildet sich eine Holliday-Struktur aus, die über die Gene 2 und 3 hinweg wandert. Zeichnen Sie die Schritte der Rekombination und lösen Sie die Holliday-Struktur auf zwei Arten auf.} 
\url{https://www.youtube.com/watch?v=vLuWg1IxoJ4}
\subsubsection*{Doppelstrangbruch}
Ein Doppelstrangbruch (DSB) tritt an den Gene 2 und 3 auf.
\subsubsection*{Strang-Austausch}
Die Einzelstränge von Molekül 1 und Molekül 2 paaren sich an den komplementären Stellen.
\subsubsection*{Holliday-Struktur bilden}
Es bildet sich eine Kreuzstruktur, bei der die Stränge der beiden Moleküle über die Gene 2 und 3 hinweg verknüpft sind.
\subsubsection*{Auflösung 1: vertikal}
Bei der vertikalen Schneidung werden die DNA-Stränge entlang der langen Achse der Holliday-Struktur geschnitten. \newline 
Die Stränge werden so geschnitten, dass Gene 2 und 3 von Molekül 1 an Gene 2 und 3 von Molekül 2 getauscht werden.
\subsubsection*{Auflösung 2: horizontal}
Bei der horizontalen Schneidung werden die Stränge entlang der kurzen Achse der Holliday-Struktur geschnitten. \newline
Die Stränge werden so geschnitten, dass Gene 2 und 3 von Molekül 1 und 2 bleiben, aber in entgegengesetzter Orientierung. \newline


\subsection*{Bei welchen Genen kann Genkonversion auftreten?}
Genkonversion kann bei Genen auftreten, die eine hohe Homologie aufweisen, wie in Genfamilien oder bei repetitiven DNA-Elementen.

\subsection*{Als Oswald Avery zu einem Transformationsansatz mit einem Gemisch aus durch Hitze getöteten 5-Zellen und lebenden R-Zellen von Streptococcus pneumoniae folgende Substanz zusetzte unterblieb die Transformation und die mit dem Gemisch infizierte Maus überlebte.
}
\begin{todolist}
\item Lipase
\item Lipase + Protease 
\item Phosphatase + Protease
\item Phosphatase 
\item [\done] DNase
\item RNase
\item Protease
\end{todolist}
\subsection*{Bei Bakterien gibt es verschiedene Möglichkeiten der DNA-Übertragung. Der direkte Zell-Zell-Kontakt ist wesentlich bei der ...}
\begin{todolist}
\item Transformation
\item [\done] Konjugation
\item Transduktion
\item Phagen-Lyse 
\item Transkription
\item Lysogenie
\item Translation
\item DNA-Replikation
\end{todolist}
\subsection*{Sie transformieren einen Tryptophan-, Arginin- und Leucin- auxotrophen Bakterien-Stamm trp- arg- leu- mit DNA eines Arginin- und Leucin-prototrophen, aber Tryptophan-auxotrophen Stammes trp - arg + leu+. 
Welche Rekombinante(n) kann/können entstehen?}
trp - arg + leu+ ? 

\subsection*{Enzyme und Proteine der DNA-Replikation}
\begin{table}[H]
    \centering 
    \begin{tabular}{|p{4cm}|p{11cm}|}
    \hline
    \textbf{DnaA-Protein} &  Initiation der Replikation am oriC\\\hline 
    \textbf{SSB (Single-Strand Binding Protein)} & Bindet an einzelsträngige DNA  \\\hline
    Helikase &  \textbf{Löst die Basenpaarung von doppelten DNA- oder RNA-Strängen auf} \\\hline
    Primase &   \textbf{erzeugt RNA-Primer}\\\hline
    \textbf{DNA-Polymerase III.} & Synthetisiert den neuen DNA-Strang \\\hline
    \textbf{DNA-Polymerase I.} & Entfernt den RNA-Primer (5'-3'-Exonuklease) und füllt den Strang auf \\\hline
    Ligase &  \textbf{Verknüpfen von DNA-Strängen}\\\hline
    Topoisomerase & \textbf{entspannen Supercoiling} \\\hline
\end{tabular}
\end{table}
\subsection*{tRNA}
\subsubsection*{Wie lang sind tRNA-Moleküle?}
etwa 70 bis 90 Nukleotide
\subsubsection*{Welche Struktur nehmen tRNA-Moleküle ein?}
kleeblattartige Sekundärstruktur \textit{die durch intramolekulare Basenpaarungen innerhalb des Moleküls gebildet wird}
\subsubsection*{Was kann man über die Basen in tRNA sagen?}
tRNA-Moleküle bestehen nicht nur aus den vier Standard-RNA-Basen, sondern auch aus zahlreichen modifizierten Basen, die die Struktur und Funktion der tRNA optimieren. U.a. 
\begin{itemize} 
    \item Pseudouridin
    \item Inosin
    \item Dihydrouridin
    \item Methylierte Basen
\end{itemize}

\subsubsection*{Welche Funktion hat tRNA?}
Vermittler zwischen der mRNA (messenger RNA) und den Aminosäuren in der Proteinbiosynthese \textit{(Translation)}. 







\end{document}